\documentclass[oneside,a4paper,14pt]{extarticle}
\usepackage[a4paper,top=20mm,bottom=20mm,left=30mm,right=15mm]{geometry}
\usepackage[T2A]{fontenc}
\usepackage[utf8]{inputenc}
\usepackage[russian]{babel}
\usepackage{indentfirst}
\usepackage{graphicx}
\usepackage{float}
\usepackage{caption}
\captionsetup[figure]{labelsep=endash, justification=centering, name=Рисунок}
\captionsetup[table]{labelsep=endash, justification=justified, singlelinecheck=false, name=Таблица}
\usepackage{titlesec}
\usepackage{minted, fancyvrb}
\usepackage{hyperref}
\usepackage{enumitem}
\usepackage{newunicodechar}
\newunicodechar{┌}{+}
\newunicodechar{─}{-}
\newunicodechar{┐}{+}
\newunicodechar{│}{\textbar}
\newunicodechar{└}{+}
\newunicodechar{┘}{+}
\newunicodechar{═}{=}
\newunicodechar{╔}{+}
\newunicodechar{╗}{+}
\newunicodechar{╚}{+}
\newunicodechar{╝}{+}
\newunicodechar{║}{\textbar}
\newunicodechar{▂}{\textunderscore}
\newunicodechar{▃}{\textunderscore}
\newunicodechar{▄}{\textunderscore}
\newunicodechar{▅}{\textunderscore}
\newunicodechar{▆}{\textunderscore}
\newunicodechar{▇}{\textunderscore}
\newunicodechar{█}{\textunderscore}
\newunicodechar{▕}{\textbar}
\newunicodechar{▏}{\textbar}
\newunicodechar{✓}{v}
\newunicodechar{✗}{x}


\setminted{style = rainbow_dash, fontsize = \scriptsize}
% Header styles
\titleformat{\section}{\normalsize\bfseries}{\thesection}{1em}{}
\titleformat{\subsection}{\normalsize\bfseries}{\thesubsection}{1em}{}
\titleformat{\subsubsection}{\normalsize\bfseries}{\thesubsubsection}{1em}{}

% Spacing
\renewcommand\baselinestretch{1.5}

\setlength{\parindent}{1.25cm}  % длина красной строки

% List styles
\setlist[enumerate]{
  left=\parindent,
  label=\arabic*.,
  itemsep=0pt,
  topsep=5pt,
  partopsep=0pt,
  parsep=0pt
}

\setlist[itemize]{
  left=\parindent,
  label=--,
  itemsep=0pt,
  topsep=5pt,
  partopsep=0pt,
  parsep=0pt
}

% Hyperlinks
\hypersetup{
  colorlinks=true,
  linkcolor=black,
  urlcolor=blue,
  pdfborder={0 0 0},
  pdftitle={Система управления гоночными командами Formula Racing},
  pdfauthor={Черкасов А.А.}
}



\begin{document}

\newpage
\thispagestyle{empty}
\begin{center}
  МИНИСТЕРСТВО НАУКИ И ВЫСШЕГО ОБРАЗОВАНИЯ РОССИЙСКОЙ ФЕДЕРАЦИИ ФЕДЕРАЛЬНОЕ ГОСУДАРСТВЕННОЕ БЮДЖЕТНОЕ ОБРАЗОВАТЕЛЬНОЕ УЧРЕЖДЕНИЕ ВЫСШЕГО ОБРАЗОВАНИЯ\\
  «ВЯТСКИЙ ГОСУДАРСТВЕННЫЙ УНИВЕРСИТЕТ»\\
  Институт математики и информационных систем\\
  Факультет автоматики и вычислительной техники\\
  Кафедра электронных вычислительных машин
\end{center}
\vspace{10mm}

\hfill
\begin{tabular}{l}
  \footnotesize Дата сдачи на проверку:                                          \\
  \footnotesize <<\rule[-1mm]{5mm}{0.10mm}\/>>\rule[-1mm]{20mm}{0.10mm}\ 2026 г. \\
  \footnotesize Проверено:                                                       \\
  \footnotesize <<\rule[-1mm]{5mm}{0.10mm}\/>>\rule[-1mm]{20mm}{0.10mm}\ 2026 г. \\
\end{tabular}
\vfill

\begin{center}
  Разработка иерархии классов Formula Racing Management System на C++. \\
  Отчёт по лабораторной работе №1\\
  по дисциплине\\
  <<Объектно-ориентированное программирование>>\\
\end{center}
\vspace{25mm}
\noindent
\begin{tabular}{ll}
  Разработал студент гр. ИВТб-2301-05-00 & \hspace{18mm}\rule[-1mm]{30mm}{0.10mm}\,/Черкасов А. А./ \\
                                         & \hspace{25.5mm}\footnotesize(подпись)                    \\
  Преподаватель                  & \hspace{18mm}\rule[-1mm]{30mm}{0.10mm}\,/Шмакова Н. А./   \\
                                         & \hspace{25.5mm}\footnotesize(подпись)                    \\
\end{tabular}

\noindent
\begin{tabular}{lp{58mm}r}
  Работа защищена &  & \hspace{13mm}<<\rule[-1mm]{5mm}{0.10mm}\/>>\rule[-1mm]{30mm}{0.10mm}\ 2026 г.
\end{tabular}
\vfill

\begin{center}
  Киров\\
  2026
\end{center}

\newpage\thispagestyle{plain}

\section*{Цели лабораторной работы}
\begin{itemize}
  \item[$-$] изучить принципы инкапсуляции, наследования и полиморфизма в языке C++;
  \item[$-$] научиться проектировать иерархию классов с использованием абстрактных базовых классов и виртуальных функций;
  \item[$-$] освоить работу с объектами через указатели базового типа и использование векторов смарт-указателей.
\end{itemize}

\section*{Задание}
Разработать приложение <<Formula Racing Management System>>, демонстрирующее принципы ООП, со следующими требованиями:
\begin{enumerate}
  \item Создать иерархию классов: один родительский (абстрактный) и три дочерних (F1Team, F2Team, FormulaETeam).
  \item В каждом классе определить не менее двух данных-членов и двух собственных методов.
  \item Для дочерних классов реализовать не менее двух унаследованных и двух перекрытых функций.
  \item Доступ к полям реализовать только через геттеры и сеттеры.
  \item Обеспечить наличие не менее 3-х конструкторов в дочерних классах.
  \item Работу с объектами осуществлять только через переменные родительского типа (полиморфизм).
  \item Каждый класс должен находиться в отдельном файле.
  \item Продемонстрировать перегрузку операторов ($+=, ==, >, <<$).
\end{enumerate}

\clearpage
\section*{Реализация приложения}

\subsection*{Предметная область}
Система предназначена для управления гоночными командами различных серий (Formula 1, Formula 2 и Formula E). Каждая серия имеет свои особенности начисления очков, бюджетные ограничения и технические параметры, что делает их идеальными кандидатами для демонстрации наследования и полиморфизма.

\subsection*{Диаграмма классов}
\begin{figure}[H]
  \centering
  \includegraphics[width=0.85\textwidth]{pics/class_diagram.pdf}
  \caption{Диаграмма классов системы управления гоночными командами}
\end{figure}

Описание иерархии:
\begin{itemize}
  \item[$-$] \texttt{FormulaTeam} (Abstract) $-$ базовый класс, хранящий общую информацию (название, руководитель, очки, победы, подиумы). Содержит чистые виртуальные методы displayInfo(), getSeriesName(), calculatePoints();
  \item[$-$] \texttt{F1Team} $-$ дочерний класс, добавляющий параметры бюджета, позиции в чемпионате и силовой установки.
  \item[$-$] \texttt{F2Team} $-$ дочерний класс, фокусирующийся на подготовке пилотов (выпускники F1) и модели шасси.
  \item[$-$] \texttt{FormulaETeam} $-$ дочерний класс, учитывающий экологический рейтинг, емкость батарей и энергетических партнеров.
  \item[$-$] \texttt{FormulaDriver} $-$ вспомогательный класс (агрегируется в \texttt{FormulaTeam}), хранящий информацию о гонщике (имя, фамилия, номер, набранные очки).
\end{itemize}

\subsection*{Описание структур классов}

В данном разделе приведено подробное описание полей и методов разработанных классов предметной области.

\begin{table}[htbp]
\centering
\caption{Описание членов класса FormulaDriver}
\begin{tabular}{|p{5.5cm}|p{3cm}|p{2cm}|p{5cm}|}
\hline
\textbf{Название} & \textbf{Тип} & \textbf{Доступ} & \textbf{Описание} \\ \hline
\multicolumn{4}{|c|}{\textbf{Поля}} \\ \hline
firstName & string & private & Имя пилота \\ \hline
lastName & string & private & Фамилия пилота \\ \hline
racingNumber & int & private & Гоночный номер \\ \hline
points & int & private & Количество очков \\ \hline
\multicolumn{4}{|c|}{\textbf{Методы}} \\ \hline
FormulaDriver(...) & - & public & Конструкторы \\ \hline
getFirstName(), getLastName(), getFullName() & string & public & Получение имени пилота \\ \hline
getRacingNumber(), getPoints() & int & public & Получение номера и очков \\ \hline
addPoints(pts: int) & void & public & Добавление очков пилоту \\ \hline
\end{tabular}
\end{table}

\begin{table}[htbp]
\centering
\caption{Описание членов класса FormulaTeam (Абстрактный)}
\begin{tabular}{|p{5.5cm}|p{3cm}|p{2cm}|p{5cm}|}
\hline
\textbf{Название} & \textbf{Тип} & \textbf{Доступ} & \textbf{Описание} \\ \hline
\multicolumn{4}{|c|}{\textbf{Поля}} \\ \hline
teamName & string & protected & Название команды \\ \hline
principalName & string & protected & Имя руководителя \\ \hline
headquarters & string & protected & Штаб-квартира \\ \hline
championshipPoints & int & protected & Очки в чемпионате \\ \hline
raceWins, podiums & int & protected & Победы и подиумы \\ \hline
drivers & vector$<$FormulaDriver$>$ & protected & Список пилотов команды \\ \hline
\multicolumn{4}{|c|}{\textbf{Методы}} \\ \hline
FormulaTeam(...) & - & public & Конструкторы \\ \hline
Геттеры и сеттеры & разный & public & Доступ к полям команды \\ \hline
addPoints(), addWin(), addPodium() & void & public & Изменение статистики \\ \hline
getWinRate(totalRaces: int) & double & public & Вычисление процента побед \\ \hline
addDriver(driver: FormulaDriver) & void & public & Добавление пилота \\ \hline
displayInfo(), getSeriesName() & void, string & public & Чисто виртуальные методы \\ \hline
calculatePoints(position: int) & int & public & Вычисление очков (virtual) \\ \hline
displayStats(), getSummary() & void, string & public & Вывод статистики (virtual) \\ \hline
operator+=, operator++, operator==, $>$, $<$, $<<$ & разный & public & Перегрузка операторов \\ \hline
\end{tabular}
\end{table}

\begin{table}[htbp]
\centering
\caption{Описание членов класса F1Team}
\begin{tabular}{|p{5.5cm}|p{3cm}|p{2cm}|p{5cm}|}
\hline
\textbf{Название} & \textbf{Тип} & \textbf{Доступ} & \textbf{Описание} \\ \hline
\multicolumn{4}{|c|}{\textbf{Поля}} \\ \hline
powerUnit & string & private & Поставщик двигателя \\ \hline
budgetCapMln & double & private & Бюджет в млн USD \\ \hline
constructorPos & int & private & Позиция в кубке \\ \hline
\multicolumn{4}{|c|}{\textbf{Методы}} \\ \hline
F1Team(...) & - & public & Конструкторы \\ \hline
Геттеры и сеттеры & разный & public & Доступ к полям \\ \hline
calculateBudgetAllocation(...) & double & public & Расчет части бюджета \\ \hline
getPowerUnitStatus() & string & public & Статус силовой установки \\ \hline
Переопределенные методы & разный & public & displayInfo, getSeriesName и т.д. \\ \hline
\end{tabular}
\end{table}

\begin{table}[htbp]
\centering
\caption{Описание членов класса F2Team}
\begin{tabular}{|p{5.5cm}|p{3cm}|p{2cm}|p{5cm}|}
\hline
\textbf{Название} & \textbf{Тип} & \textbf{Доступ} & \textbf{Описание} \\ \hline
\multicolumn{4}{|c|}{\textbf{Поля}} \\ \hline
chassisModel & string & private & Модель шасси \\ \hline
f1Graduates & int & private & Выпускники в F1 \\ \hline
isFeederSeries & bool & private & Статус молодежной серии \\ \hline
\multicolumn{4}{|c|}{\textbf{Методы}} \\ \hline
F2Team(...) & - & public & Конструкторы \\ \hline
Геттеры и сеттеры & разный & public & Доступ к полям \\ \hline
addGraduate() & void & public & Увеличение числа выпускников \\ \hline
getFeederStatus() & string & public & Текстовый статус серии \\ \hline
Переопределенные методы & разный & public & displayInfo, getSeriesName и т.д. \\ \hline
\end{tabular}
\end{table}

\begin{table}[htbp]
\centering
\caption{Описание членов класса FormulaETeam}
\begin{tabular}{|p{5.5cm}|p{3cm}|p{2cm}|p{5cm}|}
\hline
\textbf{Название} & \textbf{Тип} & \textbf{Доступ} & \textbf{Описание} \\ \hline
\multicolumn{4}{|c|}{\textbf{Поля}} \\ \hline
sustainabilityScore & int & private & Экологический рейтинг \\ \hline
batteryCapacityKwh & double & private & Емкость батареи (кВтч) \\ \hline
energyPartner & string & private & Энергетический партнер \\ \hline
\multicolumn{4}{|c|}{\textbf{Методы}} \\ \hline
FormulaETeam(...) & - & public & Конструкторы \\ \hline
Геттеры и сеттеры & разный & public & Доступ к полям \\ \hline
updateSustainabilityScore(...) & void & public & Обновление эко-рейтинга \\ \hline
getBatteryEfficiencyRating() & string & public & Оценка батареи \\ \hline
Переопределенные методы & разный & public & displayInfo, getSeriesName и т.д. \\ \hline
\end{tabular}
\end{table}


\subsection*{Особенности реализации}

\subsubsection*{Абстрактный базовый класс}
Базовый класс \texttt{FormulaTeam} является абстрактным, так как содержит чистые виртуальные функции. Это гарантирует, что мы не сможем создать экземпляр общего типа <<Команда>>, а будем работать только с конкретными реализациями.

\begin{minted}{cpp}
virtual void displayInfo() const = 0;
virtual std::string getSeriesName() const = 0;
virtual int calculatePoints(int position) const = 0;
\end{minted}

\subsubsection*{Полиморфизм, векторы и смарт-указатели}
Для автоматического управления временем жизни объектов в куче используются умные указатели \texttt{std::shared\_ptr} и \texttt{std::unique\_ptr}. В системе используется контейнер \texttt{std::vector} для хранения полиморфных указателей на базовый класс \texttt{FormulaTeam}. Объекты дочерних классов создаются динамически и помещаются в этот вектор (upcasting), что позволяет хранить и обрабатывать их единообразно:

\begin{minted}{cpp}
std::vector<std::unique_ptr<FormulaTeam>> poly_teams;

// Добавление объектов дочерних классов в вектор
poly_teams.push_back(std::make_unique<F1Team>("McLaren", "Stella", "UK", "Mercedes", 135.0, 2, 291, 5, 12));
poly_teams.push_back(std::make_unique<F2Team>("Prema", "Rosin", "Italy", "Dallara F2", 0, true, 89, 2, 5));

// Чтение из вектора и вызов виртуальных методов (полиморфизм)
for (size_t i = 0; i < poly_teams.size(); ++i) {
    std::string series = poly_teams[i]->getSeriesName(); // Вызов переопределенного метода
    std::string summary = poly_teams[i]->getSummary();
}
\end{minted}

\subsubsection*{Приведение типов (static\_cast и dynamic\_cast)}
Для вызова специфичных методов дочерних классов, которые отсутствуют в базовом классе \texttt{FormulaTeam}, применяются операции нисходящего приведения типов (downcasting). 

В приложении используется \texttt{dynamic\_cast} для безопасного приведения типов во время выполнения. Он использует механизм RTTI (Run-Time Type Information) и возвращает \texttt{nullptr}, если преобразование невозможно:

\begin{minted}{cpp}
if (auto *f1 = dynamic_cast<F1Team *>(poly_teams[0].get())) {
    // Успешно: объект действительно является F1Team
    std::string puStatus = f1->getPowerUnitStatus();
} else {
    // Ошибка приведения типов, объект относится к другой серии
}
\end{minted}

Также для случаев, когда тип объекта точно известен программе на этапе компиляции, применяется быстрое статическое приведение типов без дополнительных проверок с помощью \texttt{static\_cast}:

\begin{minted}{cpp}
if (poly_teams[0]->getSeriesName() == "Formula 1") {
    F1Team *f1 = static_cast<F1Team *>(poly_teams[0].get());
    std::string puStatus = f1->getPowerUnitStatus();
}
\end{minted}


\subsubsection*{Перегрузка операторов}
Реализована перегрузка операторов для удобной работы с объектами:
\begin{itemize}
  \item \texttt{operator+=} $-$ добавление очков;
  \item \texttt{operator==} $-$ сравнение команд по названию;
  \item \texttt{operator\textgreater} $-$ сравнение по количеству очков (используется для сортировки);
  \item \texttt{operator\textless\textless} $-$ потоковый вывод краткой информации.
\end{itemize}

\section*{Вывод}
В ходе выполнения лабораторной работы была разработана иерархия классов для управления гоночными командами. Были успешно продемонстрированы ключевые принципы ООП: инкапсуляция через приватные поля и методы доступа, наследование функционала базового класса и полиморфизм при обработке контейнера объектов разных типов через указатель базового класса. Использование раздельной компиляции (каждый класс в своем файле) делает проект масштабируемым и легкочитаемым.

\newpage

\section*{Приложение А1. Базовый класс FormulaTeam.h}
\inputminted{cpp}{code/src/core/FormulaTeam.h}

\section*{Приложение А2. Базовый класс FormulaTeam.cpp}
\inputminted{cpp}{code/src/core/FormulaTeam.cpp}

\section*{Приложение А3. Класс F1Team.h}
\inputminted{cpp}{code/src/core/F1Team.h}

\section*{Приложение А4. Класс F1Team.cpp}
\inputminted{cpp}{code/src/core/F1Team.cpp}

\section*{Приложение А5. Класс F2Team.h}
\inputminted{cpp}{code/src/core/F2Team.h}

\section*{Приложение А6. Класс F2Team.cpp}
\inputminted{cpp}{code/src/core/F2Team.cpp}

\section*{Приложение А7. Класс FormulaETeam.h}
\inputminted{cpp}{code/src/core/FormulaETeam.h}

\section*{Приложение А8. Класс FormulaETeam.cpp}
\inputminted{cpp}{code/src/core/FormulaETeam.cpp}

\section*{Приложение А9. Класс FormulaDriver.h}
\inputminted{cpp}{code/src/core/FormulaDriver.h}

\section*{Приложение А10. Класс FormulaDriver.cpp}
\inputminted{cpp}{code/src/core/FormulaDriver.cpp}

\section*{Приложение А11. Главный файл main.cpp}
\inputminted{cpp}{code/src/main.cpp}

\section*{Приложение А12. Фабрика команд TeamFactory.h}
\inputminted{cpp}{code/src/core/TeamFactory.h}

\section*{Приложение А13. Фабрика команд TeamFactory.cpp}
\inputminted{cpp}{code/src/core/TeamFactory.cpp}

\section*{Приложение А14. Класс интерфейса UI.h}
\inputminted{cpp}{code/src/ui/UI.h}

\section*{Приложение А15. Класс интерфейса UI.cpp}
\inputminted{cpp}{code/src/ui/UI.cpp}

\end{document}
